\section{Long Context、多模型协作与在线学习}

前面讨论多模态理解和生成，本章进入未来两年的两个 GPT-4 时刻。张祥雨认为，long context 很重要，但今天直接把超长上下文塞进一个 Transformer 里有问题：注意力会涣散，上下文越长，相似信息越多，模型越容易被干扰。真正的智能不是把所有东西无损记住，而是能压缩、遗忘、翻书、转移注意力。

\begin{figure}[H]
\centering
\includegraphics[width=0.96\textwidth]{long-context-multi-model.png}
\caption{Long Context 与多模型协作：不用把整本书塞进一个上下文，而是全局规划加局部执行。自制概念图，依据 01:46:56--02:07:09 对谈内容整理。}
\end{figure}

\begin{knowledgebox}{读图：全局规划和局部执行可以分开}
图中 Global Planner 保留整体印象，Retrieve/Turn page 负责翻书，Local Worker 做局部推理，Feedback 回传结果，Context Shift 切换情景。这比把所有 token 都塞进同一个上下文更接近人的注意力机制。
\end{knowledgebox}

\subsection{为什么架构不是本质}

本节解释“Linear Transformer 不本质”。张祥雨不是否认线性注意力或 RNN-like 架构的价值，而是说：如果目标只是通过架构硬扛超长上下文，可能没有触到问题核心。长上下文的核心是如何压缩、如何检索、如何切换情景、如何让多个模型或模块协作，而不只是把注意力复杂度从平方降下来。

\begin{warningbox}{不要把大海捞针当成智能本身}
为了通过 retrieve 测试，模型可能学到“什么都不能忘”的 bias。但人类读书后写总结时，会记住结构，再按需要翻书找句子。真正有用的 long context 系统，应当允许遗忘、摘要、翻页、局部复查和上下文切换。
\end{warningbox}

\begin{knowledgebox}{术语消化：long context 的两类问题}
\noindent\begin{tabularx}{\textwidth}{p{0.20\textwidth}p{0.35\textwidth}X}
\toprule
问题 & 含义 & 本期中的判断 \\
\midrule
计算复杂度 & 注意力随长度增长带来成本问题 & 重要，但不是唯一问题。 \\
智能建模 & 如何压缩、遗忘、转移注意力和局部复查 & 更接近真正的长程智能。 \\
多模型协作 & 用 planner 和 worker 分工处理上下文 & 可减少干扰，并让系统可训练。 \\
情景隔离 & 做第二题时不必背着第一题全部上下文 & 避免长上下文污染当前推理。 \\
\bottomrule
\end{tabularx}
\end{knowledgebox}

\subsection{在线学习与自主学习}

如果 long context 是一个潜在 GPT-4 时刻，另一个就是在线学习或自主学习。张祥雨把它和 Agent 联系起来：今天很多 Agent 是应用层工具编排，但更根本的 Agent 应当能在环境中持续行动、接收反馈、改进策略。这里的 scaling 不再只是 model scaling 或 data scaling，而是 environmental scaling：能不能提供足够多、足够真实、足够可验证的环境反馈。

\begin{figure}[H]
\centering
\includegraphics[width=0.96\textwidth]{online-learning-agent.png}
\caption{在线学习与自主学习：下一个 GPT-4 时刻可能来自环境反馈和持续自我改进。自制概念图，依据 02:08:30--02:25:00 对谈内容整理。}
\end{figure}

\begin{knowledgebox}{读图：环境反馈是新的 scaling 对象}
图中 Environment、Agent、Feedback、RL/IRL 和 Online Learning 连成闭环。模型不再只从离线语料中学习，而是在任务、工具、世界和人类反馈中持续更新。难点是环境规模、反馈质量和安全边界。
\end{knowledgebox}

\begin{importantbox}{课堂提示：Agent 应用和自主学习 Agent 不是一回事}
今天许多 Agent 产品是工具编排和工作流自动化；张祥雨讨论的更根本 Agent，是能在环境反馈中持续改进的学习系统。前者是应用形态，后者是训练范式和智能机制。
\end{importantbox}

\subsection{世界模型与机器人抢跑}

访谈最后把多模态和机器人连接起来。人没有视觉生成器官，但人有世界模型：我们可以在脑中预测空间、物理和行动后果。机器学习系统是否必须通过生成式训练获得世界模型，仍是开放问题；但如果机器人要在物理世界行动，它不能绕过视觉理解、动作反馈和世界模型。

\begin{figure}[H]
\centering
\includegraphics[width=0.94\textwidth]{world-model-robotics-gap.png}
\caption{世界模型与机器人抢跑：人没有生成器官，但有世界模型；机器人不能跳过视觉理解。自制概念图，依据 02:25:00--02:28:44 对谈内容整理。}
\end{figure}

\begin{knowledgebox}{读图：机器人不能只押硬件身体}
左边 World Model 负责空间推理、物理常识和未来预测；右边 Robotics 需要视觉、动作、反馈和本体共同过线。机器人行业若跳过视觉推理和世界模型，只靠本体和遥操作，很容易在智能上抢跑。
\end{knowledgebox}

\subsection{本章小结}

未来两个关键突破方向是 long context 的真实智能建模，以及在线/自主学习的环境反馈闭环。前者要求模型学会压缩、翻书和情景切换；后者要求模型在行动中持续改进。

